\documentclass[12pt,a4paper]{article}
\usepackage[UTF8]{ctex}
\usepackage{amsmath,amssymb,amsthm}
\usepackage{geometry}

\geometry{margin=2.5cm}

\title{叉积的反对称性}
\author{第8章补充说明}
\date{}

\begin{document}

\maketitle

\section{问题提出}

\textbf{问题}：什么是叉积的反对称性？

\textbf{符号说明}：
\begin{itemize}
\item $a, b, c \in \mathbb{R}^3$：三维向量
\item $a \times b$：向量$a$和$b$的叉积
\item $[a]$或$[a]_{\times}$：向量$a$的反对称矩阵
\end{itemize}

\section{叉积的反对称性}

\subsection{基本定义}

\textbf{叉积的反对称性}是指：
\begin{equation}
a \times b = -(b \times a) \tag{1}
\end{equation}

或者等价地：
\begin{equation}
a \times b = -b \times a
\end{equation}

\subsection{几何意义}

叉积的反对称性反映了叉积的方向性：
\begin{itemize}
\item $a \times b$的方向由右手定则确定
\item $b \times a$的方向与$a \times b$相反
\item 两者的大小相同：$\|a \times b\| = \|b \times a\| = \|a\| \|b\| \sin\theta$
\end{itemize}

\subsection{直接证明}

从叉积的定义出发：
\begin{equation}
a \times b = \begin{bmatrix}
a_y b_z - a_z b_y \\
a_z b_x - a_x b_z \\
a_x b_y - a_y b_x
\end{bmatrix}
\end{equation}

计算$b \times a$：
\begin{equation}
b \times a = \begin{bmatrix}
b_y a_z - b_z a_y \\
b_z a_x - b_x a_z \\
b_x a_y - b_y a_x
\end{bmatrix} = -\begin{bmatrix}
a_y b_z - a_z b_y \\
a_z b_x - a_x b_z \\
a_x b_y - a_y b_x
\end{bmatrix} = -(a \times b)
\end{equation}

\textbf{因此}：$a \times b = -b \times a$。

\section{反对称矩阵的反对称性}

\subsection{反对称矩阵的定义}

对于向量$a = [a_x, a_y, a_z]^T$，其反对称矩阵定义为：
\begin{equation}
[a] = [a]_{\times} = \begin{bmatrix}
0 & -a_z & a_y \\
a_z & 0 & -a_x \\
-a_y & a_x & 0
\end{bmatrix}
\end{equation}

\subsection{反对称矩阵的反对称性}

\textbf{反对称矩阵的反对称性}是指：
\begin{equation}
[a]^T = -[a] \tag{2}
\end{equation}

\textbf{证明}：

计算$[a]^T$：
\begin{equation}
[a]^T = \begin{bmatrix}
0 & a_z & -a_y \\
-a_z & 0 & a_x \\
a_y & -a_x & 0
\end{bmatrix} = -\begin{bmatrix}
0 & -a_z & a_y \\
a_z & 0 & -a_x \\
-a_y & a_x & 0
\end{bmatrix} = -[a]
\end{equation}

\textbf{因此}：$[a]^T = -[a]$。

\subsection{一般反对称矩阵的定义}

一个矩阵$A$是反对称的（skew-symmetric），如果：
\begin{equation}
A^T = -A
\end{equation}

这意味着：
\begin{itemize}
\item 对角线元素为零：$A_{ii} = 0$
\item 非对角线元素满足：$A_{ij} = -A_{ji}$
\end{itemize}

\section{叉积与反对称矩阵的关系}

\subsection{基本关系}

叉积可以用反对称矩阵表示为：
\begin{equation}
a \times b = [a] b
\end{equation}

\textbf{验证}：
\begin{align}
[a] b &= \begin{bmatrix}
0 & -a_z & a_y \\
a_z & 0 & -a_x \\
-a_y & a_x & 0
\end{bmatrix} \begin{bmatrix}
b_x \\ b_y \\ b_z
\end{bmatrix} \\
&= \begin{bmatrix}
-a_z b_y + a_y b_z \\
a_z b_x - a_x b_z \\
-a_y b_x + a_x b_y
\end{bmatrix} \\
&= \begin{bmatrix}
a_y b_z - a_z b_y \\
a_z b_x - a_x b_z \\
a_x b_y - a_y b_x
\end{bmatrix} \\
&= a \times b
\end{align}

\subsection{反对称性的等价形式}

由于$a \times b = [a] b$和$a \times b = -b \times a$，我们有：
\begin{equation}
[a] b = -[b] a \tag{3}
\end{equation}

\textbf{证明}：
\begin{align}
[a] b &= a \times b \\
&= -(b \times a) \\
&= -[b] a
\end{align}

\textbf{因此}：$[a] b = -[b] a$。

\section{反对称性的应用}

\subsection{应用1：叉积的交换}

利用反对称性，我们可以交换叉积中的向量：
\begin{align}
a \times b &= -b \times a \\
[a] b &= -[b] a
\end{align}

\subsection{应用2：标量三重积}

对于标量三重积$(a \times b) \cdot c$：
\begin{align}
(a \times b) \cdot c &= ([a] b) \cdot c \\
&= b^T [a]^T c \\
&= -b^T [a] c \quad \text{（使用$[a]^T = -[a]$）} \\
&= -b \cdot (a \times c)
\end{align}

\subsection{应用3：向量三重积}

对于向量三重积$a \times (b \times c)$：
\begin{align}
a \times (b \times c) &= [a] (b \times c) \\
&= [a] [b] c \\
&= -[a] [c] b \quad \text{（使用$[b] c = -[c] b$）}
\end{align}

\section{反对称性的不同层面}

\subsection{层面1：叉积运算的反对称性}

\begin{equation}
a \times b = -b \times a
\end{equation}

这是叉积运算本身的反对称性。

\subsection{层面2：反对称矩阵的反对称性}

\begin{equation}
[a]^T = -[a]
\end{equation}

这是矩阵的反对称性。

\subsection{层面3：反对称矩阵乘积的反对称性}

\begin{equation}
[a] b = -[b] a
\end{equation}

这是反对称矩阵与向量乘积的反对称性。

\section{详细推导}

\subsection{推导1：从叉积定义出发}

\textbf{叉积的定义}：
\begin{equation}
a \times b = \begin{bmatrix}
a_y b_z - a_z b_y \\
a_z b_x - a_x b_z \\
a_x b_y - a_y b_x
\end{bmatrix}
\end{equation}

\textbf{交换向量}：
\begin{equation}
b \times a = \begin{bmatrix}
b_y a_z - b_z a_y \\
b_z a_x - b_x a_z \\
b_x a_y - b_y a_x
\end{bmatrix}
\end{equation}

\textbf{比较}：
\begin{align}
b \times a &= \begin{bmatrix}
b_y a_z - b_z a_y \\
b_z a_x - b_x a_z \\
b_x a_y - b_y a_x
\end{bmatrix} \\
&= -\begin{bmatrix}
a_y b_z - a_z b_y \\
a_z b_x - a_x b_z \\
a_x b_y - a_y b_x
\end{bmatrix} \\
&= -(a \times b)
\end{align}

\textbf{因此}：$a \times b = -b \times a$。

\subsection{推导2：从反对称矩阵出发}

\textbf{反对称矩阵的定义}：
\begin{equation}
[a] = \begin{bmatrix}
0 & -a_z & a_y \\
a_z & 0 & -a_x \\
-a_y & a_x & 0
\end{bmatrix}
\end{equation}

\textbf{转置}：
\begin{equation}
[a]^T = \begin{bmatrix}
0 & a_z & -a_y \\
-a_z & 0 & a_x \\
a_y & -a_x & 0
\end{bmatrix}
\end{equation}

\textbf{比较}：
\begin{align}
[a]^T &= \begin{bmatrix}
0 & a_z & -a_y \\
-a_z & 0 & a_x \\
a_y & -a_x & 0
\end{bmatrix} \\
&= -\begin{bmatrix}
0 & -a_z & a_y \\
a_z & 0 & -a_x \\
-a_y & a_x & 0
\end{bmatrix} \\
&= -[a]
\end{align}

\textbf{因此}：$[a]^T = -[a]$。

\subsection{推导3：从反对称矩阵乘积出发}

\textbf{已知}：
\begin{align}
a \times b &= [a] b \\
a \times b &= -b \times a
\end{align}

\textbf{因此}：
\begin{align}
[a] b &= a \times b \\
&= -(b \times a) \\
&= -[b] a
\end{align}

\textbf{因此}：$[a] b = -[b] a$。

\section{重要恒等式}

\subsection{恒等式1：叉积的反对称性}

\begin{equation}
a \times b = -b \times a
\end{equation}

\subsection{恒等式2：反对称矩阵的转置}

\begin{equation}
[a]^T = -[a]
\end{equation}

\subsection{恒等式3：反对称矩阵乘积的交换}

\begin{equation}
[a] b = -[b] a
\end{equation}

\subsection{恒等式4：标量三重积的反对称性}

\begin{equation}
(a \times b) \cdot c = -(b \times a) \cdot c
\end{equation}

\subsection{恒等式5：向量三重积}

\begin{equation}
a \times (b \times c) = (a \cdot c) b - (a \cdot b) c
\end{equation}

\section{应用示例}

\subsection{示例1：速度合成}

在速度合成中，我们经常遇到：
\begin{equation}
\omega \times p = -p \times \omega
\end{equation}

\textbf{使用反对称矩阵}：
\begin{equation}
[\omega] p = -[p] \omega
\end{equation}

\subsection{示例2：角速度变换}

在角速度变换中：
\begin{equation}
R [\omega] R^T = [R \omega]
\end{equation}

\textbf{使用反对称性}：
\begin{align}
R [\omega] R^T &= R [\omega] R^T \\
&= -R [\omega]^T R^T \quad \text{（使用$[\omega]^T = -[\omega]$）} \\
&= -([\omega] R)^T R^T
\end{align}

\subsection{示例3：动力学方程}

在动力学方程中，科里奥利项：
\begin{equation}
\omega \times (I \omega) = [\omega] I \omega = -[I \omega] \omega
\end{equation}

\section{几何直观}

\subsection{叉积的几何意义}

叉积$a \times b$产生一个垂直于$a$和$b$的向量：
\begin{itemize}
\item 方向：由右手定则确定
\item 大小：$\|a \times b\| = \|a\| \|b\| \sin\theta$
\end{itemize}

\textbf{反对称性}：交换$a$和$b$会改变方向，但保持大小不变。

\subsection{反对称矩阵的几何意义}

反对称矩阵$[a]$表示绕轴$a$的旋转：
\begin{itemize}
\item $[a] b = a \times b$表示向量$b$绕轴$a$旋转90度后的方向
\item 反对称性$[a]^T = -[a]$表示旋转方向的相反
\end{itemize}

\section{总结}

\subsection{主要结论}

\begin{enumerate}
\item \textbf{叉积的反对称性}：
\begin{equation}
a \times b = -b \times a
\end{equation}

\item \textbf{反对称矩阵的反对称性}：
\begin{equation}
[a]^T = -[a]
\end{equation}

\item \textbf{反对称矩阵乘积的反对称性}：
\begin{equation}
[a] b = -[b] a
\end{equation}
\end{enumerate}

\subsection{关键要点}

\begin{itemize}
\item 叉积的反对称性反映了叉积的方向性
\item 反对称矩阵的反对称性是矩阵的结构性质
\item 两者通过关系$a \times b = [a] b$相互联系
\item 反对称性在机器人学中广泛应用于速度合成和动力学方程
\end{itemize}

\subsection{应用场景}

\begin{itemize}
\item 速度合成公式的推导
\item 角速度与线速度的关系
\item 动力学方程中的科里奥利项
\item 反对称矩阵的相似变换
\item 约束力的计算
\end{itemize}

\end{document}

